\documentclass[10pt,twocolumn]{article}
\usepackage[a4paper,margin=15mm]{geometry}
\usepackage{amsmath,amssymb}
\usepackage{mathrsfs}
\usepackage{amsthm}
\theoremstyle{definition}
\newtheorem{theorem}{Theorem}
\newtheorem{proposition}{Proposition}
\newtheorem{corollary}{Corollary}
\newtheorem{remark}{Remark}
\newtheorem{definition}{Definition}
\usepackage{graphicx}
\usepackage{booktabs,array,calc}
\usepackage[font=small,labelfont=bf]{caption}
\usepackage[colorlinks=true]{hyperref}
\usepackage{etoolbox}
\providecommand{\tightlist}{\setlength{\itemsep}{0pt}\setlength{\parskip}{0pt}}
\setcounter{secnumdepth}{-1}
\emergencystretch=3em
\allowdisplaybreaks
\AtBeginDocument{%
  \BeforeBeginEnvironment{equation}{\small}%
  \BeforeBeginEnvironment{equation*}{\small}%
  \BeforeBeginEnvironment{align}{\small}%
  \BeforeBeginEnvironment{align*}{\small}%
  \BeforeBeginEnvironment{gather}{\small}%
  \BeforeBeginEnvironment{gather*}{\small}%
}

\begin{document}
\title{The Policy-Class Lattice: Which Restriction\\ Destroyed Viability?}
\author{Amin Abaee\\[0.35em]
{\small Independent Researcher}\\[0.55em]
{\small\href{https://orcid.org/0000-0002-0019-1842}{ORCID: 0000-0002-0019-1842}}\\[0.3em]
{\small\href{mailto:amin\_abaee@ut.ac.ir}{amin\_abaee@ut.ac.ir}}}
\date{September 23, 2026}
\maketitle

\begin{abstract}
Which restriction destroyed viability --- the codex, the reading, the
instrument set, or the belief? The programme's papers separate policy
classes case by case; this paper, the D7 item of the generalization
catalogue, organizes them as a lattice with a computable diagnostic. For
policy classes on the two-patch working system, viability kernels are
monotone under class inclusion, the chain institutional $\subseteq$
observation-based $\subseteq$ full-information holds setwise with exact
counts $24\subseteq 26\subseteq 28$ of $36$ safe pairs and is strict at
both links, and the certainty-equivalence separation mechanism of the
successor's trap (corrected-class drift $\ge\tfrac{21}{100}$) is
re-verified as the off-lattice separator. The diagnostic contribution is
\emph{distance to viability}: for the audited nonviable belief, the
minimum sensor refinement that restores viability is exactly one split,
the minimum action-set expansion is exactly one instrument (which restores
viability under the unchanged reading, enlarging the viable family to all
$36$ pairs), and the minimum prior shrinking is exactly one state. Every
number is machine-verified in exact arithmetic (6/6 check families).
\end{abstract}

\section{Introduction and relation to the programme}

The completed items distinguish perfect-information policies,
observation-based policies, certainty-equivalence policies, and
institutionally restricted policies as separate mechanisms. The source
catalogue proposes the unified reading adopted here: a parameterized class
$\Pi$ with kernels $\mathrm{Viab}_\Pi(\mathcal V)$, monotone under
inclusion, forming a hierarchy --- and, more consequentially for
governance, a diagnostic question: which restriction is responsible for the
loss? The present paper instantiates that programme on the audited working
system, computing the inclusions and the distances exactly. No completed
item is modified; the stochastic selector's probabilistic class and the
successor's five layers are cited where they bound the reading.

\section{Policy classes and the chain}

A policy class $\Pi$ is a family of observation-based laws (measurable in
the reading, acting per cell under the successor's read-then-act
semantics). Its kernel is the set of beliefs from which some law of the
class maintains the floors indefinitely.

\begin{proposition}[Monotonicity]
$\Pi_1\subseteq\Pi_2\ \Longrightarrow\
\mathrm{Viab}_{\Pi_1}(\mathcal V)\subseteq
\mathrm{Viab}_{\Pi_2}(\mathcal V)$.
\end{proposition}
\begin{proof}
A witness law for $B$ under $\Pi_1$ is a law of $\Pi_2$.
\end{proof}

\begin{proposition}[The audited chain is strict]
\label{prop:chain}
On the two-patch pair family: the institutional class (aggregate reading
with a no-repeat codex) sustains $24$ of $36$ safe pairs; the per-cell
aggregate class sustains $26$; the full-information class sustains $28$.
Each inclusion is setwise and strict, with separating beliefs
$\{(1,2),(3,1)\}$ at the lower link and $\{(1,2),(2,1)\}$ at the upper.
On this family the certainty-equivalence class coincides with the
per-cell class (laws act on the reading only); the successor's trap
corrected-class drift $\ge\tfrac{21}{100}$ on $s\in[1,2]$ --- is the
mechanism that separates the two off this family.
\end{proposition}
\begin{proof}
The inclusions follow from monotonicity (each restricted class is contained
in the next). The counts and the two separating beliefs are computed
exactly by the verification script; the drift identity
$g(s+\tfrac1{10})-g(s)=\tfrac s5+\tfrac1{100}\ge\tfrac{21}{100}$ for
$s\ge1$ is exact.
\end{proof}

\section{Distance to viability}

For a nonviable belief and a family of restorative moves, the
\emph{distance} is the minimum number of moves that restores viability.

\begin{proposition}[Distances of the audited belief]
\label{prop:dist}
For $B^\ast=\{(1,2),(2,1)\}$ under the aggregate reading (nonviable):
(i) the minimum sensor refinement restoring viability is exactly one cell
split --- the unique one-split partition (the two singletons) is viable,
and zero splits is the nonviable status quo; (ii) the minimum action-set
expansion is exactly one instrument --- adding $u=3$ (protect both patches)
restores viability of $B^\ast$ under the \emph{unchanged} aggregate
reading and enlarges the aggregate viable family to all $36$ pairs
(setwise, strict); (iii) the minimum prior shrinking is exactly one state
--- either deletion yields a viable belief.
\end{proposition}
\begin{proof}
(i) A two-state belief admits exactly one partition into two nonempty
cells; the script verifies the singleton partition viable and the
zero-split (merged) verdict nonviable. (ii) With $u=3$ the per-cell
action sets of $B^\ast$'s single aggregate cell gain a jointly safe
instrument; the recursion confirms viability, and the family count
$26\to36$ is computed over all pairs. (iii) Both singletons are viable
under the unchanged reading; the script verifies both and the nonviability
of the whole.
\end{proof}

\section{Verified record (6/6 exact check families)}

E1 the setwise chain with counts $24\subseteq26\subseteq28$; E2 strictness
at both links plus the re-verified CE drift bound; E3 action-expansion
monotonicity with counts $26<36$ and a strict instance; E4 the
sensor-refinement distance $=1$ certified by exhaustion of split counts;
E5 the action-expansion distance $=1$ with the restored belief; E6 the
prior-shrinking distance $=1$ with both deletions viable.

\section{What is not claimed}

The taxonomy is instantiated on one working system at one belief family
(pairs): no general computational theory of distances is claimed, no
continuous-state version, and no claim that one instrument of protection
is a realistic policy lever --- it is the minimal formal expansion.
Institutional restrictions are modeled only as action codices; procedural
and enforcement constraints belong to the programme's decentralized-
observation companion. The stochastic selector's chance constraints are
not revisited. The certainty-equivalence identification is stated for this
family only, with the trap cited for the general separation.


\section*{Declarations}
\textbf{Funding.} None.\quad \textbf{Competing interests.} None.\quad
\textbf{Data availability.} No external data were used.\quad
\textbf{Code availability.} The verification script
\texttt{policy\_class\_lattice\_v1\_verify.py} is publicly available at
\url{https://github.com/MIKEAA2020/general-sustainability} (folder
\texttt{arena agent 1/paper rewrites/latex}); it reproduces all six check
families verbatim.\quad
\textbf{AI declaration.} GLM (Z.ai), Qwen (Alibaba Cloud) and DeepSeek AI
assisted with drafting and iterative review.

\end{document}
